\documentclass[twoside]{article}
\usepackage{xeCJK}
\setCJKmainfont{SimSun}[ItalicFont=KaiTi, BoldFont=SimHei]
\setCJKsansfont{SimHei}
\setCJKmonofont{SimSun}
\renewcommand{\figurename}{图}
\renewcommand{\tablename}{表}
\renewcommand{\refname}{参考文献}

% \usepackage{aistats2021}
% If your paper is accepted, change the options for the package
% aistats2021 as follows:
%
\usepackage[accepted]{aistats2021}
%
% This option will print headings for the title of your paper and
% headings for the authors names, plus a copyright note at the end of
% the first column of the first page.

% If you set papersize explicitly, activate the following three lines:
\special{papersize = 8.5in, 11in}
\setlength{\pdfpageheight}{11in}
\setlength{\pdfpagewidth}{8.5in}

% If you use natbib package, activate the following three lines:
\usepackage[round]{natbib}
% \renewcommand{\bibname}{References}
% \renewcommand{\bibsection}{\subsubsection*{\bibname}}

% If you use BibTeX in apalike style, activate the following line:
\bibliographystyle{apalike}

\usepackage[utf8]{inputenc} % allow utf-8 input
% \usepackage[T1]{fontenc}    % use 8-bit T1 fonts
\usepackage{hyperref}       % hyperlinks
\usepackage{url}            % simple URL typesetting
\usepackage{booktabs}       % professional-quality tables
\usepackage{amsfonts}       % blackboard math symbols
\usepackage{nicefrac}       % compact symbols for 1/2, etc.
\usepackage{microtype}      % microtypography
\usepackage{graphicx}
\usepackage{multirow}       % multirow cells for tables
\usepackage{bm}             % bold math
\usepackage{bbm}            % blackboard bold math
\usepackage{algorithm}
% \usepackage{algorithmicx}
\usepackage[noend]{algpseudocode}
\usepackage{amsthm}
\renewcommand{\proofname}{证明}
\usepackage{amsmath}

\graphicspath{ {./figures/} }

\newtheorem{lemma}{引理}
\newtheorem{theorem}{定理}
\newtheorem*{theorem*}{定理}

% Better comment layout
\newcommand{\MYCOMMENT}[1]{\State \textit{\# #1}}
% \newcommand{\MYCOMMENT}[1]{\hspace{1em} \textit{// #1}}
\newcommand{\eqnref}[1]{Equation~\ref{#1}}

% Mathematical shorthands
\newcommand{\Sample}{\textsc{\small REGSAMP}}
\DeclareMathOperator{\Tr}{Tr}
\DeclareMathOperator{\argmax}{argmax}
\DeclareMathOperator{\Lang}{Lang}
\newcommand{\uMPS}{\mbox{u-MPS}}
\newcommand{\bigO}{\mc{O}}
\newcommand{\Z}{\mc{Z}}
\newcommand{\E}{\mc{E}}
\newcommand{\EL}{\mc{E}^{\scriptstyle \ell}}
\newcommand{\ER}{\mc{E}^{\scriptstyle r}}
\newcommand{\QL}{Q_{\ell}}
\newcommand{\QR}{Q_{r}}
\newcommand{\sR}{|s|_R}
\newcommand{\sRone}{|s_1|_{R_1}}
\newcommand{\sRtwo}{|s_2|_{R_2}}
\newcommand{\Pt}{\tilde{P}}
\newcommand{\R}{\mathbb{R}}
\newcommand{\starchar}{0}
\newcommand{\alphamat}{\alpha\alpha^T}
\newcommand{\omegamat}{\omega\omega^T}
\newcommand{\ntrain}{N_{\mathrm{train}}}
\newcommand{\field}[1]{\R^{#1}}
\newcommand{\mc}[1]{\mathcal{#1}}
\newcommand{\seq}[2]{{#1}_{1}, {#1}_{2}, \ldots, {#1}_{#2}}

\begin{document}
\runningtitle{Tensor Networks for Probabilistic Sequence Modeling}

\twocolumn[

\aistatstitle{概率序列建模的张量网络\\（Tensor Networks for Probabilistic Sequence Modeling）}

\aistatsauthor{Jacob Miller \And Guillaume Rabusseau \And John Terilla}
\aistatsaddress{Mila and DIRO \\
Université de Montréal \\
\href{mailto:jmjacobmiller@gmail.com}{\texttt{jmjacobmiller@gmail.com}} \And
CCAI chair - Mila and DIRO \\
Université de Montréal \\
\href{mailto:grabus@iro.umontreal.ca}{\texttt{grabus@iro.umontreal.ca}} \And
CUNY and Tunnel \\
City University of New York \\
\href{mailto:jterilla@gc.cuny.edu}{\texttt{jterilla@gc.cuny.edu}} } ]

\begin{abstract}

Tensor networks are a powerful modeling framework developed for computational many-body physics, which have only recently been applied within machine learning. In this work we utilize a uniform matrix product state (u-MPS) model for probabilistic modeling of sequence data. We first show that u-MPS enable sequence-level parallelism, with length-$n$ sequences able to be evaluated in depth $O(\log n)$. We then introduce a novel generative algorithm giving trained u-MPS the ability to efficiently sample from a wide variety of conditional distributions, each one defined by a regular expression. Special cases of this algorithm correspond to autoregressive and fill-in-the-blank sampling, but more complex regular expressions permit the generation of richly structured data in a manner that has no direct analogue in neural generative models. Experiments on sequence modeling with synthetic and real text data show u-MPS outperforming a variety of baselines and effectively generalizing their predictions in the presence of limited data.

\end{abstract}

